% =========================================================
% Chapter: Introduction
% Purpose:
% - Establishes the context and motivation for deepfake detection
% - Summarizes background, threat model, research gaps, objectives, and contributions
% Navigation:
% - Sections below are ordered from broad background to concrete thesis objectives
% Tips:
% - Keep citations in sync with references/references.tex
% - If adding new sections, also update the Chapter overview comments here
% =========================================================
\chapter{Introduction}

The rapid progress of generative modeling has transformed the production and dissemination of digital media. Techniques grounded in Generative Adversarial Networks (GANs), variational autoencoders (VAEs), and, more recently, diffusion models have enabled the synthesis of images, audio, and videos with increasingly photo-realistic quality. While these technologies have legitimate applications in creative industries, accessibility and fidelity have also accelerated the misuse of synthetic media for disinformation, fraud, and identity abuse. As a result, reliable deepfake detection has emerged as a critical problem in multimedia forensics and trustworthy AI. This thesis advances the state of deepfake detection along two complementary axes. First, it studies a video-only hybrid architecture that jointly reasons over spatial artifacts and temporal inconsistencies in facial dynamics. Second, it extends detection to the multimodal setting and explicitly models audio–visual synchrony, leveraging viseme–phoneme temporal correlations that are characteristic of natural speech but remain difficult to faithfully reproduce in synthetic pipelines.

\section{Background}
Early deepfake generation pipelines commonly exhibited visible and audible artifacts, including boundary inconsistencies, temporal flicker, and vocoder-induced spectral anomalies. Over time, GAN-based methods improved high-frequency consistency and global realism, while VAEs facilitated face-swapping via efficient latent representations. Diffusion-based models further reduced perceptual artifacts by iteratively denoising structured content, achieving global coherence that approaches photographic quality \cite{kaur2024deepfake,rossler2019faceforensics}. In parallel, audio synthesis has benefitted from high-fidelity text-to-speech and voice conversion, powered by learned vocoders and prosody models that can mimic speaker identity and speaking style. The co-evolution of visual and acoustic synthesis has lead to hybrid deepfakes that coordinate both modalities and thus evade unimodal detectors.

\section{The Deepfake Threat}
The societal impact of high-quality deepfakes spans several domains. In news and political contexts, manipulated videos can fabricate evidence or statements to influence public opinion; in financial fraud and social engineering, deepfakes enable convincing impersonation for unauthorized transactions or access. The proliferation of user-friendly deepfake toolchains enables non-experts to generate deceptive content at scale. This development places pressure on forensic systems to maintain detection performance under compression, domain shifts, and adverse conditions typical of in-the-wild media circulation. Ensuring equitable performance across demographic groups also rises in importance as deployment scales.

\section{Limitations of Existing Detection Methods}
Existing detectors frequently specialize in a single modality and exhibit limited generalization to unseen manipulations and domains \cite{mubarak2023survey}. Visual detectors often rely on spatial artifacts or single-frame inconsistencies, achieving strong performance on curated datasets but suffering performance drops when exposed to new forgery families or distribution shifts \cite{rossler2019faceforensics}. Temporal models that process frame sequences can capture motion cues, yet may still miss multi-modal inconsistencies. Audio-only detectors identify vocoder or prosody artifacts but fail when audio streams are authentic while the video is manipulated, or vice versa. Furthermore, many detectors exhibit sensitivity to compression and resolution degradation, which are commonplace in social media. Finally, differential performance across demographic subgroups has been reported, raising concerns about fairness and reliability in real-world use.

\section{Research Gaps}
Analysis of the literature reveals several gaps that motivate the contributions of this thesis:
\begin{enumerate}[leftmargin=*,label=(\alph*)]
    \item \textbf{Spatial–temporal integration in video-only detection:} Many visual detectors emphasize spatial cues and underutilize temporal information, limiting robustness to artifacts that manifest across frames.
    \item \textbf{Multimodal synchrony as a forensic cue:} Hybrid deepfakes increasingly coordinate both audio and video streams; nevertheless, subtle viseme–phoneme misalignments persist as telltale cues. Few systems explicitly learn cross-modal temporal correlations at fine time scales.
    \item \textbf{Generalization across datasets and manipulations:} Open-world performance often degrades due to overfitting to dataset-specific artifacts. Aggregated, diverse training and augmentation strategies remain under-explored in a unified framework.
    \item \textbf{Robustness to compression and resolution:} Many methods are brittle under practical degradations, calling for architectures and training protocols that remain reliable under lossy conditions.
    \item \textbf{Fairness across demographics:} Differential error rates across gender and race groups necessitate stratified evaluation and data design choices that mitigate bias.
\end{enumerate}

\section{Thesis Objectives}
The thesis pursues the following objectives:
\begin{enumerate}[leftmargin=*,label=\arabic*.]
    \item Design and evaluate a hybrid video-only architecture that integrates spatial feature extraction with bidirectional temporal modeling and attention to capture subtle visual forgeries.
    \item Construct a synchronicity-aware audio–visual detector that explicitly models cross-modal temporal relationships using Bi-GRUs and cross-modal attention.
    \item Develop unified data preprocessing and augmentation pipelines that promote generalization across multiple benchmarks and manipulation types.
    \item Establish evaluation protocols that measure accuracy, AUC-ROC, precision–recall trade-offs, and robustness under compression and resolution changes.
    \item Perform fairness-aware analysis with stratified subgroup evaluation to quantify and reduce demographic performance gaps.
\end{enumerate}

\section{Contributions}
This thesis makes the following contributions:
\begin{enumerate}[leftmargin=*,label=\arabic*.]
    \item \textbf{Hybrid video-only detection:} An architecture that combines a ResNet50 backbone for spatial encoding with Bidirectional GRUs for temporal dynamics and Multi-Head Attention for sequence focus. Trained on an aggregated dataset (FaceForensics++, Celeb-DF, DFDC), it demonstrates strong generalization and high accuracy/AUC on mixed-domain evaluations.
    \item \textbf{Synchronicity-aware multimodal detection:} A dual-stream framework that processes video and audio with specialized encoders and per-stream Bi-GRUs, fused via 8-head cross-modal attention to detect short-duration viseme–phoneme misalignments. On a consolidated FakeAVCeleb and AV-Deepfake1M++ corpus, the model achieves 92.5\% Accuracy and 0.94 AUC-ROC, outperforming unimodal and naive fusion baselines.
    \item \textbf{Robustness and fairness evaluation:} Comprehensive analyses under lossy compression and stratified demographic assessments that inform responsible deployment considerations.
    \item \textbf{A unified experimental and preprocessing protocol:} Consistent pipelines for frame sampling, spectrogram extraction, normalization, and augmentation that are reused across visual-only and multimodal settings to facilitate reproducibility and fair comparisons.
\end{enumerate}

\section{Thesis Organisation}
This thesis is organized into eight chapters as follows: Chapter~1 introduces the problem context, research gaps, objectives, and contributions. Chapter~2 presents the Motivation, articulating practical significance, key challenges, and research questions. Chapter~3 surveys the Literature, covering visual, audio, and multimodal detection, datasets, generalization, and fairness. Chapter~4 details Methodology~I (Hybrid Visual), including data, architecture, and training protocol. Chapter~5 presents Methodology~II (Synchronicity-Aware Multimodal), including cross-modal attention and alignment. Chapter~6 reports Experiments and Results for both phases, including baselines, ablations, robustness, and fairness analyses. Chapter~7 provides the Discussion, mapping findings to gaps, deployment implications, limitations, and ethics. Chapter~8 concludes and outlines Future Work, including long-range temporal modeling and cross-lingual extensions.

% Citations referenced in this chapter are defined in the References section and will be numbered according to the established order.